
        \documentclass[fleqn]{article}      
        \usepackage{amsmath}
        \usepackage{fontspec}
        \usepackage{kotex} % 한국어 지원  

        \begin{document}      
        쉬움: \\\\  
1. 정적분의 기호는 무엇인가요?  
   a) ∫  
   b) ∑  
   c) Δ  
   d) ∏  
\\\\  
2. 다음 중 부정적분의 결과로 올바른 것은?  
   a) ∫x dx = (1/2)x² + C  
   b) ∫x dx = x² + C  
   c) ∫x dx = 2x + C  
   d) ∫x dx = (1/3)x³ + C  
\\\\  
3. ∫2x dx를 계산하면 어떤 결과가 나오나요?  
   a) x² + C  
   b) 2x² + C  
   c) x² + 2C  
   d) 2x² + 2C  
\\\\  
4. 다음 중 정적분의 의미는 무엇인가요?  
   a) 함수의 미분  
   b) 함수의 넓이  
   c) 함수의 기울기  
   d) 함수의 값  
\\\\  
5. ∫1 dx의 결과는 무엇인가요?  
   a) x + C  
   b) 1 + C  
   c) 0 + C  
   d) C  
\\\\  

중간: \\\\  
1. ∫(3x² + 2x + 1) dx의 결과는 무엇인가요?  
   a) x³ + x² + x + C  
   b) x³ + x² + C  
   c) x³ + x² + x + 1  
   d) 3x³/3 + 2x²/2 + x + C  
\\\\  
2. ∫[0, 1] (2x) dx의 값을 구하세요.  
\\\\  
3. ∫(sin x) dx의 결과는 무엇인가요?  
   a) -cos x + C  
   b) cos x + C  
   c) sin x + C  
   d) -sin x + C  
\\\\  
4. ∫(x³) dx를 계산하세요.  
\\\\  
5. ∫[1, 3] (x²) dx의 값을 구하세요.  
\\\\  

어려움: \\\\  
1. ∫(e^x) dx의 결과는 무엇인가요?  
   a) e^x + C  
   b) e^x  
   c) xe^x + C  
   d) ln(x) + C  
\\\\  
2. ∫(x * cos x) dx를 부분적분을 사용하여 구하세요.  
\\\\  
3. ∫[0, π] (sin x) dx의 값을 구하세요.  
\\\\  
4. ∫(1/x) dx의 결과는 무엇인가요?  
   a) ln(x) + C  
   b) 1/x + C  
   c) e^x + C  
   d) x + C  
\\\\  
5. ∫(x² * e^x) dx를 부분적분을 사용하여 구하세요.  
\\\\  

      
        \end{document} 
    